\documentclass[aspectratio=169,10pt]{beamer}
\makeatletter\def\input@path{{../../estilo/}{../estilo/}{estilo/}}\makeatother
\usepackage{devmobile}
\modulo[Kotlin]{07}{Estado, Navegação e Animações}{Estado em Compose, recomposição, state hoisting, efeitos colaterais, ciclo de vida, Navigation 3 e animações}

\begin{document}
\capa

\begin{frame}{Objetivos do módulo}
  \begin{conceito}[Ao final deste módulo você será capaz de]
    \begin{itemize}
      \item Criar estado com \cod{remember} e \cod{mutableStateOf} e entender a recomposição.
      \item Preservar estado em rotações com \cod{rememberSaveable}.
      \item Aplicar \emph{state hoisting} para criar componentes reutilizáveis e testáveis.
      \item Usar efeitos colaterais com segurança (\cod{LaunchedEffect}, \cod{DisposableEffect}\ldots).
      \item Explicar o ciclo de vida da Activity e coletar fluxos respeitando-o.
      \item Implementar navegação entre telas com \textbf{Navigation 3}.
      \item Animar mudanças de estado com as APIs de animação do Compose.
    \end{itemize}
  \end{conceito}
\end{frame}

\sumario

% =============================================================================
%  Fundamentos compartilhados --- Módulo 07: estado, ciclo de vida, navegação
%  e movimento
%  Incluído pelas trilhas Kotlin e Flutter com % =============================================================================
%  Fundamentos compartilhados --- Módulo 07: estado, ciclo de vida, navegação
%  e movimento
%  Incluído pelas trilhas Kotlin e Flutter com % =============================================================================
%  Fundamentos compartilhados --- Módulo 07: estado, ciclo de vida, navegação
%  e movimento
%  Incluído pelas trilhas Kotlin e Flutter com \input{../fundamentos/f07-estado-navegacao}
% =============================================================================
\section{Fundamentos: estado, ciclo de vida e navegação}

\begin{frame}{Estado: o que a tela precisa lembrar}
  \begin{columns}[T]
    \begin{column}{0.56\textwidth}
      \begin{definicao}[Estado]
        Todo valor que pode mudar com o tempo e influencia o que a tela mostra. A interface é uma função dele: \textbf{UI = f(estado)}.
      \end{definicao}
      \scriptsize
      \begin{tabularx}{\linewidth}{@{}l L l@{}}
        \toprule
        \cabecalho{Tipo} & \cabecalho{Exemplo no LabMonitor} & \cabecalho{Onde mora} \\
        \midrule
        Efêmero & painel expandido, rolagem, aba ativa & no componente \\
        Da tela & bancadas carregadas, filtro, erro & no dono da tela \\
        Do app & usuário logado, tema escolhido & em um serviço compartilhado \\
        Persistente & reservas, preferências & banco ou disco (Módulo 10) \\
        \bottomrule
      \end{tabularx}
    \end{column}
    \begin{column}{0.42\textwidth}
      \begin{definicao}[Estado derivado]
        Valor calculado a partir de outros, como o total de bancadas livres. Não se guarda: calcula-se, para nunca ficar desatualizado.
      \end{definicao}
      \begin{conceito}[Fonte única da verdade]
        Cada dado tem um único dono que pode alterá-lo; os demais apenas leem. Duas cópias editáveis do mesmo dado acabam divergindo.
      \end{conceito}
    \end{column}
  \end{columns}
\end{frame}

\begin{frame}{Fluxo unidirecional de dados}
  \begin{columns}[c]
    \begin{column}{0.46\textwidth}
      \centering
      \begin{tikzpicture}[every node/.style={font=\scriptsize}]
        \node[caixaambar,minimum width=2.4cm] (e) at (0,1.5) {\textbf{Estado}};
        \node[caixa,minimum width=2.4cm] (u) at (2,0) {\textbf{Interface}};
        \node[caixa,minimum width=2.4cm] (v) at (0,-1.5) {\textbf{Evento}\\ \tiny toque, digitação};
        \node[caixaescura,minimum width=2.4cm] (d) at (-2,0) {\textbf{Dono do estado}\\ \tiny aplica as regras};
        \draw[seta] (e) to[bend left=20] node[rotulo,right]{desce} (u);
        \draw[setaambar] (u) to[bend left=20] node[rotulo,right]{gera} (v);
        \draw[setaambar] (v) to[bend left=20] node[rotulo,left]{sobe} (d);
        \draw[seta] (d) to[bend left=20] node[rotulo,left]{novo valor} (e);
      \end{tikzpicture}
    \end{column}
    \begin{column}{0.52\textwidth}
      \begin{definicao}[Fluxo unidirecional (UDF)]
        O estado desce do dono até a interface; os eventos sobem da interface até o dono. \textbf{Só o dono altera o estado}.
      \end{definicao}
      \begin{conceito}[Vantagens]
        Previsível: cada mudança tem um único caminho. Fácil de depurar: registre eventos e estados. Fácil de testar: um evento entra, um estado sai.
      \end{conceito}
      \begin{historia}
        A Arquitetura Elm, o Flux do Facebook (2014) e o Redux (2015) popularizaram o modelo, hoje padrão no Compose, no Flutter e no SwiftUI.
      \end{historia}
    \end{column}
  \end{columns}
\end{frame}

\begin{frame}{A tela como máquina de estados}
  \begin{columns}[c]
    \begin{column}{0.5\textwidth}
      \centering
      \begin{tikzpicture}[every node/.style={font=\scriptsize}]
        \node[caixaescura,minimum width=2cm] (c) at (0,0) {Carregando};
        \node[caixa,minimum width=1.8cm,fill=verdeok!12,draw=verdeok] (s) at (3,1.1) {Sucesso};
        \node[caixa,minimum width=1.8cm] (v) at (3,0) {Vazio};
        \node[caixa,minimum width=1.8cm,fill=vermelhoerro!10,draw=vermelhoerro] (x) at (3,-1.1) {Erro};
        \draw[seta] (-1.6,0) node[rotulo,left]{abrir tela} -- (c);
        \draw[seta] (c) -- node[rotulo,above,sloped]{dados} (s.west);
        \draw[seta] (c) -- node[rotulo,above]{lista vazia} (v);
        \draw[seta] (c) -- node[rotulo,below,sloped]{falha} (x.west);
        \draw[setaambar,rounded corners] (x.south) -- ++(0,-0.4) -| node[rotulo,pos=0.25,below]{tentar de novo} (c.south);
        \draw[setaambar,rounded corners] (s.north) -- ++(0,0.4) -| node[rotulo,pos=0.25,above]{puxar para atualizar} (c.north);
      \end{tikzpicture}
    \end{column}
    \begin{column}{0.48\textwidth}
      \begin{definicao}[Máquina de estados finitos]
        Conjunto fechado de estados e de transições permitidas, disparadas por eventos. A tela está sempre em \textbf{exatamente um} estado.
      \end{definicao}
      \begin{atencao}[Estados impossíveis]
        Três booleanos (\cod{carregando}, \cod{erro}, \cod{vazio}) geram 8 combinações, a maioria sem sentido. Uma classe selada com quatro casos admite só os quatro válidos.
      \end{atencao}
      \begin{dica}
        Kotlin e Dart 3 têm classes seladas: o compilador obriga o \cod{when} ou o \cod{switch} a tratar todos os casos.
      \end{dica}
    \end{column}
  \end{columns}
\end{frame}

\begin{frame}{O que sobrevive a quê}
  \scriptsize
  \newcommand{\vsim}{\textcolor{verdeok}{\textbf{sim}}}
  \newcommand{\vnao}{\textcolor{vermelhoerro}{não}}
  \begin{tabularx}{\textwidth}{@{}L *{4}{>{\centering\arraybackslash}p{1.95cm}}@{}}
    \toprule
    \cabecalho{Onde o valor está} & \cabecalho{nova recomposição} & \cabecalho{rotação da tela} & \cabecalho{morte do processo} & \cabecalho{app fechado pelo usuário} \\
    \midrule
    variável local dentro da função de interface & \vnao & \vnao & \vnao & \vnao \\
    \cod{remember} (Compose) & \vsim & \vnao & \vnao & \vnao \\
    \cod{rememberSaveable} (Compose) & \vsim & \vsim & \vsim & \vnao \\
    \cod{ViewModel} (Android) & \vsim & \vsim & \vnao$^1$ & \vnao \\
    \cod{State} de um \cod{StatefulWidget} (Flutter) & \vsim & \vsim$^2$ & \vnao & \vnao \\
    \cod{RestorationMixin} (Flutter) & \vsim & \vsim & \vsim & \vnao \\
    banco de dados, arquivo, preferências (Módulo 10) & \vsim & \vsim & \vsim & \vsim \\
    \bottomrule
  \end{tabularx}\\[1pt]
  {\tiny $^1$ sim, com \cod{SavedStateHandle}. \quad $^2$ o Flutter não recria a tela ao girar.}
  \normalsize
  \begin{columns}[T]
    \begin{column}{0.49\textwidth}
      \begin{definicao}[Morte do processo]
        Em segundo plano, o sistema pode encerrar o app para liberar memória. Ao voltar, o usuário espera a tela como a deixou: o sistema a recria e devolve o pequeno estado salvo.
      \end{definicao}
    \end{column}
    \begin{column}{0.49\textwidth}
      \begin{dica}[Como testar]
        Ative ``Não manter atividades'' nas opções do desenvolvedor do Android, ou rode \cod{adb shell am kill <pacote>} com o app em segundo plano.
      \end{dica}
    \end{column}
  \end{columns}
\end{frame}

\begin{frame}{Ciclo de vida e prioridade dos processos}
  \begin{columns}[c]
    \begin{column}{0.44\textwidth}
      \centering
      \begin{tikzpicture}[node distance=1.5mm,every node/.style={font=\scriptsize}]
        \node[caixaescura,minimum width=3.6cm] (a) {primeiro plano (em uso)};
        \node[caixapetroleo,minimum width=3.6cm,below=of a] (b) {visível, mas encoberto};
        \node[caixa,minimum width=3.6cm,below=of b] (c) {serviço em execução};
        \node[caixa,minimum width=3.6cm,below=of c,fill=cinzaclaro] (d) {segundo plano (em cache)};
        \draw[seta] ([xshift=2mm]a.north east) -- node[rotulo,right,text width=1.4cm]{mais importante} ([xshift=2mm]d.south east);
        \node[rotulo,below=2mm of d,text width=3.8cm] {com pouca memória, o sistema encerra de baixo para cima};
      \end{tikzpicture}
    \end{column}
    \begin{column}{0.54\textwidth}
      \begin{definicao}[Ciclo de vida]
        Sequência de estados por que passam o app e cada tela, controlada pelo sistema: criado, visível, em primeiro plano, em segundo plano, destruído. Cada transição é avisada por \emph{callbacks}.
      \end{definicao}
      \begin{conceito}[O que fazer nas transições]
        Ao sair do primeiro plano: pare câmera, sensores e animações e salve rascunhos. Ao voltar: retome e atualize os dados.
      \end{conceito}
      \begin{pordentro}[Por que o sistema encerra apps]
        A memória do celular é limitada, e o sistema prefere encerrar apps em segundo plano a deixar o aparelho lento: no Android, o \emph{low memory killer}; no iOS, o \emph{jetsam}.
      \end{pordentro}
    \end{column}
  \end{columns}
\end{frame}

\begin{frame}{Navegação: grafo de destinos e pilha de retorno}
  \begin{columns}[c]
    \begin{column}{0.48\textwidth}
      \centering
      \begin{tikzpicture}[node distance=4mm and 3mm,every node/.style={font=\tiny}]
        \node[caixa,minimum width=1.4cm,font=\tiny] (l) {Login};
        \node[caixaescura,minimum width=1.4cm,font=\tiny,right=of l] (p) {Painel};
        \node[caixa,minimum width=1.4cm,font=\tiny,below=of p] (e) {Equipamentos};
        \node[caixa,minimum width=1.4cm,font=\tiny,right=of p] (a) {Ajustes};
        \node[caixa,minimum width=1.4cm,font=\tiny,below=of e] (d) {Detalhe};
        \node[caixaambar,minimum width=1.4cm,font=\tiny,below=of d] (r) {Reserva};
        \draw[seta] (l) -- (p); \draw[seta] (p) -- (e); \draw[seta] (p) -- (a); \draw[seta] (e) -- (d); \draw[seta] (d) -- (r);
        \node[rotulo,below=1mm of r] {grafo};
        \node[caixa,minimum width=1.9cm,font=\tiny] (s1) at (4.6,-2.6) {Painel};
        \node[caixa,minimum width=1.9cm,font=\tiny,above=0.5mm of s1] (s2) {Equipamentos};
        \node[caixaambarforte,minimum width=1.9cm,font=\tiny,above=0.5mm of s2] (s3) {Detalhe (topo)};
        \node[pontilhado,fit=(s1)(s2)(s3)] (pil) {};
        \node[rotulo,below=1mm of pil] {pilha};
      \end{tikzpicture}
    \end{column}
    \begin{column}{0.5\textwidth}
      \begin{definicao}[Grafo e pilha]
        Destinos (telas) são \textbf{nós}; as ligações permitidas, \textbf{arestas}. A pilha de retorno é LIFO: navegar empilha, voltar desempilha.
      \end{definicao}
      \begin{conceito}[Princípios]
        Destino inicial fixo; o voltar nunca surpreende; cada aba da barra inferior tem a própria pilha.
      \end{conceito}
      \begin{atencao}[Voltar preditivo]
        Com \cod{targetSdk} 36 no Android 16, o gesto mostra a tela anterior antes de o dedo soltar e \cod{onBackPressed} não é mais chamado: use \cod{BackHandler} ou \cod{PopScope}.
      \end{atencao}
    \end{column}
  \end{columns}
\end{frame}

\begin{frame}{Navegação declarativa e links profundos}
  \begin{columns}[T]
    \begin{column}{0.5\textwidth}
      \begin{definicao}[Imperativa \emph{versus} declarativa]
        Imperativa: comandos \cod{push} e \cod{pop}. Declarativa: a pilha é um \textbf{estado} (uma lista de destinos) e o que aparece é função dela, como na interface (Navigation 3, go\_router).
      \end{definicao}
      \begin{definicao}[Link profundo (\emph{deep link})]
        Endereço que abre o app direto em um destino, como \cod{https://labmonitor.app/reserva/42}. O app monta uma pilha sintética (Painel, Reserva) para o voltar fazer sentido.
      \end{definicao}
      \begin{atencao}
        Parâmetros de um link vêm de fora: valide-os (a reserva existe? este usuário pode vê-la?).
      \end{atencao}
    \end{column}
    \begin{column}{0.48\textwidth}
      \centering
      \begin{tikzpicture}[node distance=3mm,every node/.style={font=\tiny}]
        \node[caixa,minimum width=3.6cm,font=\tiny] (o) {e-mail, notificação ou QR Code};
        \node[caixapetroleo,minimum width=3.6cm,font=\tiny,below=of o] (s) {sistema confere o domínio};
        \node[caixaescura,minimum width=3.6cm,font=\tiny,below=of s] (r) {roteador do app interpreta};
        \node[caixaambar,minimum width=3.6cm,font=\tiny,below=of r] (p) {pilha: Painel $\rightarrow$ Reserva 42};
        \draw[seta] (o) -- (s); \draw[seta] (s) -- (r); \draw[seta] (r) -- (p);
      \end{tikzpicture}
      \raggedright
      \begin{pordentro}[Links verificados]
        \emph{App Links} (Android, arquivo \cod{assetlinks.json}) e \emph{Universal Links} (iOS, \cod{apple-app-site-association}) provam que o site e o app têm o mesmo dono.
      \end{pordentro}
    \end{column}
  \end{columns}
\end{frame}

\begin{frame}{Movimento: interpolação, curvas e molas}
  \begin{columns}[T]
    \begin{column}{0.44\textwidth}
      \centering
      \begin{tikzpicture}[x=3.4cm,y=1.9cm,baseline=(current bounding box.north)]
        \draw[->,cinzamedio] (0,0) -- (1.08,0) node[rotulo,right]{t};
        \draw[->,cinzamedio] (0,0) -- (0,1.35) node[rotulo,above]{progresso};
        \draw[cinzaclaro,dashed] (0,1) -- (1,1);
        \draw[cinzamedio,line width=0.8pt] plot[domain=0:1,samples=2] (\x,\x);
        \draw[azulpetroleo,line width=1pt] plot[domain=0:1,samples=40] (\x,{1-(1-\x)^3});
        \draw[ambar,line width=1pt] plot[domain=0:1,samples=80] (\x,{1-exp(-5*\x)*cos(deg(10*\x))});
        \node[rotulo,text=cinzamedio] at (0.82,0.62) {linear};
        \node[rotulo,text=azulpetroleo] at (0.3,0.88) {desacelera};
        \node[rotulo,text=ambar] at (0.42,1.3) {mola (passa e volta)};
      \end{tikzpicture}
      \raggedright
      \begin{definicao}[Interpolação]
        A cada quadro: \cod{valor = início + (fim - início)} $\times$ \cod{curva(t)}, com \cod{t} indo de 0 a 1 ao longo da animação.
      \end{definicao}
    \end{column}
    \begin{column}{0.54\textwidth}
      \begin{conceito}[Curvas e molas]
        Curvas de aceleração (\emph{easing}) imitam objetos reais: desaceleram ao chegar. \textbf{Molas} são definidas por rigidez e amortecimento, não por duração, e podem ser interrompidas sem saltos. O Material 3 Expressive (2025) adotou molas como base.
      \end{conceito}
      \begin{dica}[Movimento com propósito]
        Anime para orientar: de onde veio a tela, o que mudou. Transições de 150 a 400\,ms mantêm o app ágil.
      \end{dica}
      \begin{atencao}[Acessibilidade]
        Respeite a opção do sistema de remover ou reduzir animações: movimento excessivo causa desconforto em algumas pessoas.
      \end{atencao}
    \end{column}
  \end{columns}
\end{frame}

% =============================================================================
\section{Fundamentos: estado, ciclo de vida e navegação}

\begin{frame}{Estado: o que a tela precisa lembrar}
  \begin{columns}[T]
    \begin{column}{0.56\textwidth}
      \begin{definicao}[Estado]
        Todo valor que pode mudar com o tempo e influencia o que a tela mostra. A interface é uma função dele: \textbf{UI = f(estado)}.
      \end{definicao}
      \scriptsize
      \begin{tabularx}{\linewidth}{@{}l L l@{}}
        \toprule
        \cabecalho{Tipo} & \cabecalho{Exemplo no LabMonitor} & \cabecalho{Onde mora} \\
        \midrule
        Efêmero & painel expandido, rolagem, aba ativa & no componente \\
        Da tela & bancadas carregadas, filtro, erro & no dono da tela \\
        Do app & usuário logado, tema escolhido & em um serviço compartilhado \\
        Persistente & reservas, preferências & banco ou disco (Módulo 10) \\
        \bottomrule
      \end{tabularx}
    \end{column}
    \begin{column}{0.42\textwidth}
      \begin{definicao}[Estado derivado]
        Valor calculado a partir de outros, como o total de bancadas livres. Não se guarda: calcula-se, para nunca ficar desatualizado.
      \end{definicao}
      \begin{conceito}[Fonte única da verdade]
        Cada dado tem um único dono que pode alterá-lo; os demais apenas leem. Duas cópias editáveis do mesmo dado acabam divergindo.
      \end{conceito}
    \end{column}
  \end{columns}
\end{frame}

\begin{frame}{Fluxo unidirecional de dados}
  \begin{columns}[c]
    \begin{column}{0.46\textwidth}
      \centering
      \begin{tikzpicture}[every node/.style={font=\scriptsize}]
        \node[caixaambar,minimum width=2.4cm] (e) at (0,1.5) {\textbf{Estado}};
        \node[caixa,minimum width=2.4cm] (u) at (2,0) {\textbf{Interface}};
        \node[caixa,minimum width=2.4cm] (v) at (0,-1.5) {\textbf{Evento}\\ \tiny toque, digitação};
        \node[caixaescura,minimum width=2.4cm] (d) at (-2,0) {\textbf{Dono do estado}\\ \tiny aplica as regras};
        \draw[seta] (e) to[bend left=20] node[rotulo,right]{desce} (u);
        \draw[setaambar] (u) to[bend left=20] node[rotulo,right]{gera} (v);
        \draw[setaambar] (v) to[bend left=20] node[rotulo,left]{sobe} (d);
        \draw[seta] (d) to[bend left=20] node[rotulo,left]{novo valor} (e);
      \end{tikzpicture}
    \end{column}
    \begin{column}{0.52\textwidth}
      \begin{definicao}[Fluxo unidirecional (UDF)]
        O estado desce do dono até a interface; os eventos sobem da interface até o dono. \textbf{Só o dono altera o estado}.
      \end{definicao}
      \begin{conceito}[Vantagens]
        Previsível: cada mudança tem um único caminho. Fácil de depurar: registre eventos e estados. Fácil de testar: um evento entra, um estado sai.
      \end{conceito}
      \begin{historia}
        A Arquitetura Elm, o Flux do Facebook (2014) e o Redux (2015) popularizaram o modelo, hoje padrão no Compose, no Flutter e no SwiftUI.
      \end{historia}
    \end{column}
  \end{columns}
\end{frame}

\begin{frame}{A tela como máquina de estados}
  \begin{columns}[c]
    \begin{column}{0.5\textwidth}
      \centering
      \begin{tikzpicture}[every node/.style={font=\scriptsize}]
        \node[caixaescura,minimum width=2cm] (c) at (0,0) {Carregando};
        \node[caixa,minimum width=1.8cm,fill=verdeok!12,draw=verdeok] (s) at (3,1.1) {Sucesso};
        \node[caixa,minimum width=1.8cm] (v) at (3,0) {Vazio};
        \node[caixa,minimum width=1.8cm,fill=vermelhoerro!10,draw=vermelhoerro] (x) at (3,-1.1) {Erro};
        \draw[seta] (-1.6,0) node[rotulo,left]{abrir tela} -- (c);
        \draw[seta] (c) -- node[rotulo,above,sloped]{dados} (s.west);
        \draw[seta] (c) -- node[rotulo,above]{lista vazia} (v);
        \draw[seta] (c) -- node[rotulo,below,sloped]{falha} (x.west);
        \draw[setaambar,rounded corners] (x.south) -- ++(0,-0.4) -| node[rotulo,pos=0.25,below]{tentar de novo} (c.south);
        \draw[setaambar,rounded corners] (s.north) -- ++(0,0.4) -| node[rotulo,pos=0.25,above]{puxar para atualizar} (c.north);
      \end{tikzpicture}
    \end{column}
    \begin{column}{0.48\textwidth}
      \begin{definicao}[Máquina de estados finitos]
        Conjunto fechado de estados e de transições permitidas, disparadas por eventos. A tela está sempre em \textbf{exatamente um} estado.
      \end{definicao}
      \begin{atencao}[Estados impossíveis]
        Três booleanos (\cod{carregando}, \cod{erro}, \cod{vazio}) geram 8 combinações, a maioria sem sentido. Uma classe selada com quatro casos admite só os quatro válidos.
      \end{atencao}
      \begin{dica}
        Kotlin e Dart 3 têm classes seladas: o compilador obriga o \cod{when} ou o \cod{switch} a tratar todos os casos.
      \end{dica}
    \end{column}
  \end{columns}
\end{frame}

\begin{frame}{O que sobrevive a quê}
  \scriptsize
  \newcommand{\vsim}{\textcolor{verdeok}{\textbf{sim}}}
  \newcommand{\vnao}{\textcolor{vermelhoerro}{não}}
  \begin{tabularx}{\textwidth}{@{}L *{4}{>{\centering\arraybackslash}p{1.95cm}}@{}}
    \toprule
    \cabecalho{Onde o valor está} & \cabecalho{nova recomposição} & \cabecalho{rotação da tela} & \cabecalho{morte do processo} & \cabecalho{app fechado pelo usuário} \\
    \midrule
    variável local dentro da função de interface & \vnao & \vnao & \vnao & \vnao \\
    \cod{remember} (Compose) & \vsim & \vnao & \vnao & \vnao \\
    \cod{rememberSaveable} (Compose) & \vsim & \vsim & \vsim & \vnao \\
    \cod{ViewModel} (Android) & \vsim & \vsim & \vnao$^1$ & \vnao \\
    \cod{State} de um \cod{StatefulWidget} (Flutter) & \vsim & \vsim$^2$ & \vnao & \vnao \\
    \cod{RestorationMixin} (Flutter) & \vsim & \vsim & \vsim & \vnao \\
    banco de dados, arquivo, preferências (Módulo 10) & \vsim & \vsim & \vsim & \vsim \\
    \bottomrule
  \end{tabularx}\\[1pt]
  {\tiny $^1$ sim, com \cod{SavedStateHandle}. \quad $^2$ o Flutter não recria a tela ao girar.}
  \normalsize
  \begin{columns}[T]
    \begin{column}{0.49\textwidth}
      \begin{definicao}[Morte do processo]
        Em segundo plano, o sistema pode encerrar o app para liberar memória. Ao voltar, o usuário espera a tela como a deixou: o sistema a recria e devolve o pequeno estado salvo.
      \end{definicao}
    \end{column}
    \begin{column}{0.49\textwidth}
      \begin{dica}[Como testar]
        Ative ``Não manter atividades'' nas opções do desenvolvedor do Android, ou rode \cod{adb shell am kill <pacote>} com o app em segundo plano.
      \end{dica}
    \end{column}
  \end{columns}
\end{frame}

\begin{frame}{Ciclo de vida e prioridade dos processos}
  \begin{columns}[c]
    \begin{column}{0.44\textwidth}
      \centering
      \begin{tikzpicture}[node distance=1.5mm,every node/.style={font=\scriptsize}]
        \node[caixaescura,minimum width=3.6cm] (a) {primeiro plano (em uso)};
        \node[caixapetroleo,minimum width=3.6cm,below=of a] (b) {visível, mas encoberto};
        \node[caixa,minimum width=3.6cm,below=of b] (c) {serviço em execução};
        \node[caixa,minimum width=3.6cm,below=of c,fill=cinzaclaro] (d) {segundo plano (em cache)};
        \draw[seta] ([xshift=2mm]a.north east) -- node[rotulo,right,text width=1.4cm]{mais importante} ([xshift=2mm]d.south east);
        \node[rotulo,below=2mm of d,text width=3.8cm] {com pouca memória, o sistema encerra de baixo para cima};
      \end{tikzpicture}
    \end{column}
    \begin{column}{0.54\textwidth}
      \begin{definicao}[Ciclo de vida]
        Sequência de estados por que passam o app e cada tela, controlada pelo sistema: criado, visível, em primeiro plano, em segundo plano, destruído. Cada transição é avisada por \emph{callbacks}.
      \end{definicao}
      \begin{conceito}[O que fazer nas transições]
        Ao sair do primeiro plano: pare câmera, sensores e animações e salve rascunhos. Ao voltar: retome e atualize os dados.
      \end{conceito}
      \begin{pordentro}[Por que o sistema encerra apps]
        A memória do celular é limitada, e o sistema prefere encerrar apps em segundo plano a deixar o aparelho lento: no Android, o \emph{low memory killer}; no iOS, o \emph{jetsam}.
      \end{pordentro}
    \end{column}
  \end{columns}
\end{frame}

\begin{frame}{Navegação: grafo de destinos e pilha de retorno}
  \begin{columns}[c]
    \begin{column}{0.48\textwidth}
      \centering
      \begin{tikzpicture}[node distance=4mm and 3mm,every node/.style={font=\tiny}]
        \node[caixa,minimum width=1.4cm,font=\tiny] (l) {Login};
        \node[caixaescura,minimum width=1.4cm,font=\tiny,right=of l] (p) {Painel};
        \node[caixa,minimum width=1.4cm,font=\tiny,below=of p] (e) {Equipamentos};
        \node[caixa,minimum width=1.4cm,font=\tiny,right=of p] (a) {Ajustes};
        \node[caixa,minimum width=1.4cm,font=\tiny,below=of e] (d) {Detalhe};
        \node[caixaambar,minimum width=1.4cm,font=\tiny,below=of d] (r) {Reserva};
        \draw[seta] (l) -- (p); \draw[seta] (p) -- (e); \draw[seta] (p) -- (a); \draw[seta] (e) -- (d); \draw[seta] (d) -- (r);
        \node[rotulo,below=1mm of r] {grafo};
        \node[caixa,minimum width=1.9cm,font=\tiny] (s1) at (4.6,-2.6) {Painel};
        \node[caixa,minimum width=1.9cm,font=\tiny,above=0.5mm of s1] (s2) {Equipamentos};
        \node[caixaambarforte,minimum width=1.9cm,font=\tiny,above=0.5mm of s2] (s3) {Detalhe (topo)};
        \node[pontilhado,fit=(s1)(s2)(s3)] (pil) {};
        \node[rotulo,below=1mm of pil] {pilha};
      \end{tikzpicture}
    \end{column}
    \begin{column}{0.5\textwidth}
      \begin{definicao}[Grafo e pilha]
        Destinos (telas) são \textbf{nós}; as ligações permitidas, \textbf{arestas}. A pilha de retorno é LIFO: navegar empilha, voltar desempilha.
      \end{definicao}
      \begin{conceito}[Princípios]
        Destino inicial fixo; o voltar nunca surpreende; cada aba da barra inferior tem a própria pilha.
      \end{conceito}
      \begin{atencao}[Voltar preditivo]
        Com \cod{targetSdk} 36 no Android 16, o gesto mostra a tela anterior antes de o dedo soltar e \cod{onBackPressed} não é mais chamado: use \cod{BackHandler} ou \cod{PopScope}.
      \end{atencao}
    \end{column}
  \end{columns}
\end{frame}

\begin{frame}{Navegação declarativa e links profundos}
  \begin{columns}[T]
    \begin{column}{0.5\textwidth}
      \begin{definicao}[Imperativa \emph{versus} declarativa]
        Imperativa: comandos \cod{push} e \cod{pop}. Declarativa: a pilha é um \textbf{estado} (uma lista de destinos) e o que aparece é função dela, como na interface (Navigation 3, go\_router).
      \end{definicao}
      \begin{definicao}[Link profundo (\emph{deep link})]
        Endereço que abre o app direto em um destino, como \cod{https://labmonitor.app/reserva/42}. O app monta uma pilha sintética (Painel, Reserva) para o voltar fazer sentido.
      \end{definicao}
      \begin{atencao}
        Parâmetros de um link vêm de fora: valide-os (a reserva existe? este usuário pode vê-la?).
      \end{atencao}
    \end{column}
    \begin{column}{0.48\textwidth}
      \centering
      \begin{tikzpicture}[node distance=3mm,every node/.style={font=\tiny}]
        \node[caixa,minimum width=3.6cm,font=\tiny] (o) {e-mail, notificação ou QR Code};
        \node[caixapetroleo,minimum width=3.6cm,font=\tiny,below=of o] (s) {sistema confere o domínio};
        \node[caixaescura,minimum width=3.6cm,font=\tiny,below=of s] (r) {roteador do app interpreta};
        \node[caixaambar,minimum width=3.6cm,font=\tiny,below=of r] (p) {pilha: Painel $\rightarrow$ Reserva 42};
        \draw[seta] (o) -- (s); \draw[seta] (s) -- (r); \draw[seta] (r) -- (p);
      \end{tikzpicture}
      \raggedright
      \begin{pordentro}[Links verificados]
        \emph{App Links} (Android, arquivo \cod{assetlinks.json}) e \emph{Universal Links} (iOS, \cod{apple-app-site-association}) provam que o site e o app têm o mesmo dono.
      \end{pordentro}
    \end{column}
  \end{columns}
\end{frame}

\begin{frame}{Movimento: interpolação, curvas e molas}
  \begin{columns}[T]
    \begin{column}{0.44\textwidth}
      \centering
      \begin{tikzpicture}[x=3.4cm,y=1.9cm,baseline=(current bounding box.north)]
        \draw[->,cinzamedio] (0,0) -- (1.08,0) node[rotulo,right]{t};
        \draw[->,cinzamedio] (0,0) -- (0,1.35) node[rotulo,above]{progresso};
        \draw[cinzaclaro,dashed] (0,1) -- (1,1);
        \draw[cinzamedio,line width=0.8pt] plot[domain=0:1,samples=2] (\x,\x);
        \draw[azulpetroleo,line width=1pt] plot[domain=0:1,samples=40] (\x,{1-(1-\x)^3});
        \draw[ambar,line width=1pt] plot[domain=0:1,samples=80] (\x,{1-exp(-5*\x)*cos(deg(10*\x))});
        \node[rotulo,text=cinzamedio] at (0.82,0.62) {linear};
        \node[rotulo,text=azulpetroleo] at (0.3,0.88) {desacelera};
        \node[rotulo,text=ambar] at (0.42,1.3) {mola (passa e volta)};
      \end{tikzpicture}
      \raggedright
      \begin{definicao}[Interpolação]
        A cada quadro: \cod{valor = início + (fim - início)} $\times$ \cod{curva(t)}, com \cod{t} indo de 0 a 1 ao longo da animação.
      \end{definicao}
    \end{column}
    \begin{column}{0.54\textwidth}
      \begin{conceito}[Curvas e molas]
        Curvas de aceleração (\emph{easing}) imitam objetos reais: desaceleram ao chegar. \textbf{Molas} são definidas por rigidez e amortecimento, não por duração, e podem ser interrompidas sem saltos. O Material 3 Expressive (2025) adotou molas como base.
      \end{conceito}
      \begin{dica}[Movimento com propósito]
        Anime para orientar: de onde veio a tela, o que mudou. Transições de 150 a 400\,ms mantêm o app ágil.
      \end{dica}
      \begin{atencao}[Acessibilidade]
        Respeite a opção do sistema de remover ou reduzir animações: movimento excessivo causa desconforto em algumas pessoas.
      \end{atencao}
    \end{column}
  \end{columns}
\end{frame}

% =============================================================================
\section{Fundamentos: estado, ciclo de vida e navegação}

\begin{frame}{Estado: o que a tela precisa lembrar}
  \begin{columns}[T]
    \begin{column}{0.56\textwidth}
      \begin{definicao}[Estado]
        Todo valor que pode mudar com o tempo e influencia o que a tela mostra. A interface é uma função dele: \textbf{UI = f(estado)}.
      \end{definicao}
      \scriptsize
      \begin{tabularx}{\linewidth}{@{}l L l@{}}
        \toprule
        \cabecalho{Tipo} & \cabecalho{Exemplo no LabMonitor} & \cabecalho{Onde mora} \\
        \midrule
        Efêmero & painel expandido, rolagem, aba ativa & no componente \\
        Da tela & bancadas carregadas, filtro, erro & no dono da tela \\
        Do app & usuário logado, tema escolhido & em um serviço compartilhado \\
        Persistente & reservas, preferências & banco ou disco (Módulo 10) \\
        \bottomrule
      \end{tabularx}
    \end{column}
    \begin{column}{0.42\textwidth}
      \begin{definicao}[Estado derivado]
        Valor calculado a partir de outros, como o total de bancadas livres. Não se guarda: calcula-se, para nunca ficar desatualizado.
      \end{definicao}
      \begin{conceito}[Fonte única da verdade]
        Cada dado tem um único dono que pode alterá-lo; os demais apenas leem. Duas cópias editáveis do mesmo dado acabam divergindo.
      \end{conceito}
    \end{column}
  \end{columns}
\end{frame}

\begin{frame}{Fluxo unidirecional de dados}
  \begin{columns}[c]
    \begin{column}{0.46\textwidth}
      \centering
      \begin{tikzpicture}[every node/.style={font=\scriptsize}]
        \node[caixaambar,minimum width=2.4cm] (e) at (0,1.5) {\textbf{Estado}};
        \node[caixa,minimum width=2.4cm] (u) at (2,0) {\textbf{Interface}};
        \node[caixa,minimum width=2.4cm] (v) at (0,-1.5) {\textbf{Evento}\\ \tiny toque, digitação};
        \node[caixaescura,minimum width=2.4cm] (d) at (-2,0) {\textbf{Dono do estado}\\ \tiny aplica as regras};
        \draw[seta] (e) to[bend left=20] node[rotulo,right]{desce} (u);
        \draw[setaambar] (u) to[bend left=20] node[rotulo,right]{gera} (v);
        \draw[setaambar] (v) to[bend left=20] node[rotulo,left]{sobe} (d);
        \draw[seta] (d) to[bend left=20] node[rotulo,left]{novo valor} (e);
      \end{tikzpicture}
    \end{column}
    \begin{column}{0.52\textwidth}
      \begin{definicao}[Fluxo unidirecional (UDF)]
        O estado desce do dono até a interface; os eventos sobem da interface até o dono. \textbf{Só o dono altera o estado}.
      \end{definicao}
      \begin{conceito}[Vantagens]
        Previsível: cada mudança tem um único caminho. Fácil de depurar: registre eventos e estados. Fácil de testar: um evento entra, um estado sai.
      \end{conceito}
      \begin{historia}
        A Arquitetura Elm, o Flux do Facebook (2014) e o Redux (2015) popularizaram o modelo, hoje padrão no Compose, no Flutter e no SwiftUI.
      \end{historia}
    \end{column}
  \end{columns}
\end{frame}

\begin{frame}{A tela como máquina de estados}
  \begin{columns}[c]
    \begin{column}{0.5\textwidth}
      \centering
      \begin{tikzpicture}[every node/.style={font=\scriptsize}]
        \node[caixaescura,minimum width=2cm] (c) at (0,0) {Carregando};
        \node[caixa,minimum width=1.8cm,fill=verdeok!12,draw=verdeok] (s) at (3,1.1) {Sucesso};
        \node[caixa,minimum width=1.8cm] (v) at (3,0) {Vazio};
        \node[caixa,minimum width=1.8cm,fill=vermelhoerro!10,draw=vermelhoerro] (x) at (3,-1.1) {Erro};
        \draw[seta] (-1.6,0) node[rotulo,left]{abrir tela} -- (c);
        \draw[seta] (c) -- node[rotulo,above,sloped]{dados} (s.west);
        \draw[seta] (c) -- node[rotulo,above]{lista vazia} (v);
        \draw[seta] (c) -- node[rotulo,below,sloped]{falha} (x.west);
        \draw[setaambar,rounded corners] (x.south) -- ++(0,-0.4) -| node[rotulo,pos=0.25,below]{tentar de novo} (c.south);
        \draw[setaambar,rounded corners] (s.north) -- ++(0,0.4) -| node[rotulo,pos=0.25,above]{puxar para atualizar} (c.north);
      \end{tikzpicture}
    \end{column}
    \begin{column}{0.48\textwidth}
      \begin{definicao}[Máquina de estados finitos]
        Conjunto fechado de estados e de transições permitidas, disparadas por eventos. A tela está sempre em \textbf{exatamente um} estado.
      \end{definicao}
      \begin{atencao}[Estados impossíveis]
        Três booleanos (\cod{carregando}, \cod{erro}, \cod{vazio}) geram 8 combinações, a maioria sem sentido. Uma classe selada com quatro casos admite só os quatro válidos.
      \end{atencao}
      \begin{dica}
        Kotlin e Dart 3 têm classes seladas: o compilador obriga o \cod{when} ou o \cod{switch} a tratar todos os casos.
      \end{dica}
    \end{column}
  \end{columns}
\end{frame}

\begin{frame}{O que sobrevive a quê}
  \scriptsize
  \newcommand{\vsim}{\textcolor{verdeok}{\textbf{sim}}}
  \newcommand{\vnao}{\textcolor{vermelhoerro}{não}}
  \begin{tabularx}{\textwidth}{@{}L *{4}{>{\centering\arraybackslash}p{1.95cm}}@{}}
    \toprule
    \cabecalho{Onde o valor está} & \cabecalho{nova recomposição} & \cabecalho{rotação da tela} & \cabecalho{morte do processo} & \cabecalho{app fechado pelo usuário} \\
    \midrule
    variável local dentro da função de interface & \vnao & \vnao & \vnao & \vnao \\
    \cod{remember} (Compose) & \vsim & \vnao & \vnao & \vnao \\
    \cod{rememberSaveable} (Compose) & \vsim & \vsim & \vsim & \vnao \\
    \cod{ViewModel} (Android) & \vsim & \vsim & \vnao$^1$ & \vnao \\
    \cod{State} de um \cod{StatefulWidget} (Flutter) & \vsim & \vsim$^2$ & \vnao & \vnao \\
    \cod{RestorationMixin} (Flutter) & \vsim & \vsim & \vsim & \vnao \\
    banco de dados, arquivo, preferências (Módulo 10) & \vsim & \vsim & \vsim & \vsim \\
    \bottomrule
  \end{tabularx}\\[1pt]
  {\tiny $^1$ sim, com \cod{SavedStateHandle}. \quad $^2$ o Flutter não recria a tela ao girar.}
  \normalsize
  \begin{columns}[T]
    \begin{column}{0.49\textwidth}
      \begin{definicao}[Morte do processo]
        Em segundo plano, o sistema pode encerrar o app para liberar memória. Ao voltar, o usuário espera a tela como a deixou: o sistema a recria e devolve o pequeno estado salvo.
      \end{definicao}
    \end{column}
    \begin{column}{0.49\textwidth}
      \begin{dica}[Como testar]
        Ative ``Não manter atividades'' nas opções do desenvolvedor do Android, ou rode \cod{adb shell am kill <pacote>} com o app em segundo plano.
      \end{dica}
    \end{column}
  \end{columns}
\end{frame}

\begin{frame}{Ciclo de vida e prioridade dos processos}
  \begin{columns}[c]
    \begin{column}{0.44\textwidth}
      \centering
      \begin{tikzpicture}[node distance=1.5mm,every node/.style={font=\scriptsize}]
        \node[caixaescura,minimum width=3.6cm] (a) {primeiro plano (em uso)};
        \node[caixapetroleo,minimum width=3.6cm,below=of a] (b) {visível, mas encoberto};
        \node[caixa,minimum width=3.6cm,below=of b] (c) {serviço em execução};
        \node[caixa,minimum width=3.6cm,below=of c,fill=cinzaclaro] (d) {segundo plano (em cache)};
        \draw[seta] ([xshift=2mm]a.north east) -- node[rotulo,right,text width=1.4cm]{mais importante} ([xshift=2mm]d.south east);
        \node[rotulo,below=2mm of d,text width=3.8cm] {com pouca memória, o sistema encerra de baixo para cima};
      \end{tikzpicture}
    \end{column}
    \begin{column}{0.54\textwidth}
      \begin{definicao}[Ciclo de vida]
        Sequência de estados por que passam o app e cada tela, controlada pelo sistema: criado, visível, em primeiro plano, em segundo plano, destruído. Cada transição é avisada por \emph{callbacks}.
      \end{definicao}
      \begin{conceito}[O que fazer nas transições]
        Ao sair do primeiro plano: pare câmera, sensores e animações e salve rascunhos. Ao voltar: retome e atualize os dados.
      \end{conceito}
      \begin{pordentro}[Por que o sistema encerra apps]
        A memória do celular é limitada, e o sistema prefere encerrar apps em segundo plano a deixar o aparelho lento: no Android, o \emph{low memory killer}; no iOS, o \emph{jetsam}.
      \end{pordentro}
    \end{column}
  \end{columns}
\end{frame}

\begin{frame}{Navegação: grafo de destinos e pilha de retorno}
  \begin{columns}[c]
    \begin{column}{0.48\textwidth}
      \centering
      \begin{tikzpicture}[node distance=4mm and 3mm,every node/.style={font=\tiny}]
        \node[caixa,minimum width=1.4cm,font=\tiny] (l) {Login};
        \node[caixaescura,minimum width=1.4cm,font=\tiny,right=of l] (p) {Painel};
        \node[caixa,minimum width=1.4cm,font=\tiny,below=of p] (e) {Equipamentos};
        \node[caixa,minimum width=1.4cm,font=\tiny,right=of p] (a) {Ajustes};
        \node[caixa,minimum width=1.4cm,font=\tiny,below=of e] (d) {Detalhe};
        \node[caixaambar,minimum width=1.4cm,font=\tiny,below=of d] (r) {Reserva};
        \draw[seta] (l) -- (p); \draw[seta] (p) -- (e); \draw[seta] (p) -- (a); \draw[seta] (e) -- (d); \draw[seta] (d) -- (r);
        \node[rotulo,below=1mm of r] {grafo};
        \node[caixa,minimum width=1.9cm,font=\tiny] (s1) at (4.6,-2.6) {Painel};
        \node[caixa,minimum width=1.9cm,font=\tiny,above=0.5mm of s1] (s2) {Equipamentos};
        \node[caixaambarforte,minimum width=1.9cm,font=\tiny,above=0.5mm of s2] (s3) {Detalhe (topo)};
        \node[pontilhado,fit=(s1)(s2)(s3)] (pil) {};
        \node[rotulo,below=1mm of pil] {pilha};
      \end{tikzpicture}
    \end{column}
    \begin{column}{0.5\textwidth}
      \begin{definicao}[Grafo e pilha]
        Destinos (telas) são \textbf{nós}; as ligações permitidas, \textbf{arestas}. A pilha de retorno é LIFO: navegar empilha, voltar desempilha.
      \end{definicao}
      \begin{conceito}[Princípios]
        Destino inicial fixo; o voltar nunca surpreende; cada aba da barra inferior tem a própria pilha.
      \end{conceito}
      \begin{atencao}[Voltar preditivo]
        Com \cod{targetSdk} 36 no Android 16, o gesto mostra a tela anterior antes de o dedo soltar e \cod{onBackPressed} não é mais chamado: use \cod{BackHandler} ou \cod{PopScope}.
      \end{atencao}
    \end{column}
  \end{columns}
\end{frame}

\begin{frame}{Navegação declarativa e links profundos}
  \begin{columns}[T]
    \begin{column}{0.5\textwidth}
      \begin{definicao}[Imperativa \emph{versus} declarativa]
        Imperativa: comandos \cod{push} e \cod{pop}. Declarativa: a pilha é um \textbf{estado} (uma lista de destinos) e o que aparece é função dela, como na interface (Navigation 3, go\_router).
      \end{definicao}
      \begin{definicao}[Link profundo (\emph{deep link})]
        Endereço que abre o app direto em um destino, como \cod{https://labmonitor.app/reserva/42}. O app monta uma pilha sintética (Painel, Reserva) para o voltar fazer sentido.
      \end{definicao}
      \begin{atencao}
        Parâmetros de um link vêm de fora: valide-os (a reserva existe? este usuário pode vê-la?).
      \end{atencao}
    \end{column}
    \begin{column}{0.48\textwidth}
      \centering
      \begin{tikzpicture}[node distance=3mm,every node/.style={font=\tiny}]
        \node[caixa,minimum width=3.6cm,font=\tiny] (o) {e-mail, notificação ou QR Code};
        \node[caixapetroleo,minimum width=3.6cm,font=\tiny,below=of o] (s) {sistema confere o domínio};
        \node[caixaescura,minimum width=3.6cm,font=\tiny,below=of s] (r) {roteador do app interpreta};
        \node[caixaambar,minimum width=3.6cm,font=\tiny,below=of r] (p) {pilha: Painel $\rightarrow$ Reserva 42};
        \draw[seta] (o) -- (s); \draw[seta] (s) -- (r); \draw[seta] (r) -- (p);
      \end{tikzpicture}
      \raggedright
      \begin{pordentro}[Links verificados]
        \emph{App Links} (Android, arquivo \cod{assetlinks.json}) e \emph{Universal Links} (iOS, \cod{apple-app-site-association}) provam que o site e o app têm o mesmo dono.
      \end{pordentro}
    \end{column}
  \end{columns}
\end{frame}

\begin{frame}{Movimento: interpolação, curvas e molas}
  \begin{columns}[T]
    \begin{column}{0.44\textwidth}
      \centering
      \begin{tikzpicture}[x=3.4cm,y=1.9cm,baseline=(current bounding box.north)]
        \draw[->,cinzamedio] (0,0) -- (1.08,0) node[rotulo,right]{t};
        \draw[->,cinzamedio] (0,0) -- (0,1.35) node[rotulo,above]{progresso};
        \draw[cinzaclaro,dashed] (0,1) -- (1,1);
        \draw[cinzamedio,line width=0.8pt] plot[domain=0:1,samples=2] (\x,\x);
        \draw[azulpetroleo,line width=1pt] plot[domain=0:1,samples=40] (\x,{1-(1-\x)^3});
        \draw[ambar,line width=1pt] plot[domain=0:1,samples=80] (\x,{1-exp(-5*\x)*cos(deg(10*\x))});
        \node[rotulo,text=cinzamedio] at (0.82,0.62) {linear};
        \node[rotulo,text=azulpetroleo] at (0.3,0.88) {desacelera};
        \node[rotulo,text=ambar] at (0.42,1.3) {mola (passa e volta)};
      \end{tikzpicture}
      \raggedright
      \begin{definicao}[Interpolação]
        A cada quadro: \cod{valor = início + (fim - início)} $\times$ \cod{curva(t)}, com \cod{t} indo de 0 a 1 ao longo da animação.
      \end{definicao}
    \end{column}
    \begin{column}{0.54\textwidth}
      \begin{conceito}[Curvas e molas]
        Curvas de aceleração (\emph{easing}) imitam objetos reais: desaceleram ao chegar. \textbf{Molas} são definidas por rigidez e amortecimento, não por duração, e podem ser interrompidas sem saltos. O Material 3 Expressive (2025) adotou molas como base.
      \end{conceito}
      \begin{dica}[Movimento com propósito]
        Anime para orientar: de onde veio a tela, o que mudou. Transições de 150 a 400\,ms mantêm o app ágil.
      \end{dica}
      \begin{atencao}[Acessibilidade]
        Respeite a opção do sistema de remover ou reduzir animações: movimento excessivo causa desconforto em algumas pessoas.
      \end{atencao}
    \end{column}
  \end{columns}
\end{frame}


% =============================================================================
\section{Estado e recomposição}

\begin{frame}[fragile]{O que é estado}
  \begin{columns}[T]
    \begin{column}{0.5\textwidth}
      \small
      \termo{Estado} é qualquer valor que pode mudar com o tempo e afeta o que aparece na tela: o texto digitado, a lista carregada, se um diálogo está aberto, a última leitura do sensor.
\begin{codigo}[title=Não funciona!]{Kotlin}
@Composable
fun Contador() {
  var cliques = 0          // variável comum
  Button(onClick = { cliques++ }) {
    Text("Cliques: $cliques")  // nunca muda
  }
}
\end{codigo}
      \begin{erro}
        Mudar uma variável comum não avisa o Compose, e a cada recomposição ela volta a 0.
      \end{erro}
    \end{column}
    \begin{column}{0.48\textwidth}
\begin{codigo}[title=Correto]{Kotlin}
@Composable
fun Contador() {
  var cliques by remember { mutableStateOf(0) }
  Button(onClick = { cliques++ }) {
    Text("Cliques: $cliques")
  }
}
\end{codigo}
      \begin{conceito}[Duas peças]
        \begin{itemize}
          \item \cod{mutableStateOf}: valor \textbf{observável}; ler registra dependência, escrever agenda recomposição.
          \item \cod{remember}: guarda o objeto entre recomposições.
        \end{itemize}
      \end{conceito}
    \end{column}
  \end{columns}
\end{frame}

\begin{frame}{Recomposição}
  \begin{columns}[c]
    \begin{column}{0.52\textwidth}
      \centering
      \begin{tikzpicture}[node distance=4mm]
        \node[caixaescura,minimum width=4.6cm] (c) {\textbf{Composição inicial}\\ executa as funções e cria a árvore};
        \node[caixaambar,minimum width=4.6cm,below=of c] (m) {\textbf{Estado muda}\\ \cod{cliques++}};
        \node[caixa,minimum width=4.6cm,below=of m] (r) {\textbf{Recomposição}\\ reexecuta \emph{só} quem leu o estado};
        \node[caixa,minimum width=4.6cm,below=of r,fill=cinzaclaro] (d) {layout e desenho atualizados};
        \draw[seta] (c) -- (m); \draw[seta] (m) -- (r); \draw[seta] (r) -- (d);
        \draw[setaambar,rounded corners] (d.east) -- ++(5mm,0) |- (m.east);
      \end{tikzpicture}
    \end{column}
    \begin{column}{0.46\textwidth}
      \begin{conceito}[Propriedades]
        \begin{itemize}
          \item É \textbf{inteligente}: funções cujos parâmetros não mudaram podem ser puladas.
          \item Pode acontecer \textbf{a cada quadro} (ex.: animações).
          \item Pode ser \textbf{cancelada} e refeita.
        \end{itemize}
      \end{conceito}
      \begin{dica}
        O \emph{Layout Inspector} mostra quantas vezes cada composable recompôs. Números muito altos indicam problema de desempenho (Módulo 14).
      \end{dica}
    \end{column}
  \end{columns}
\end{frame}

\begin{frame}[fragile]{\cod{remember} \emph{versus} \cod{rememberSaveable}}
  \begin{columns}[T]
    \begin{column}{0.52\textwidth}
\begin{codigo}[title=Sobrevivendo à rotação]{Kotlin}
// perdido ao girar a tela
var filtro by remember { mutableStateOf("") }

// salvo no Bundle: sobrevive à rotação e à
// morte do processo em segundo plano
var busca by rememberSaveable { mutableStateOf("") }

// estado derivado: só recalcula quando muda
val visiveis by remember(itens) {
  derivedStateOf { itens.filter { busca in it.descricao } }
}

// listas observáveis
val selecionados = remember { mutableStateListOf<String>() }
\end{codigo}
    \end{column}
    \begin{column}{0.46\textwidth}
      \footnotesize
      \begin{tabularx}{\linewidth}{@{}l c c@{}}
        \toprule
        & \cabecalho{recomposição} & \cabecalho{rotação} \\
        \midrule
        variável comum & \textcolor{vermelhoerro}{perde} & \textcolor{vermelhoerro}{perde} \\
        \cod{remember} & \textcolor{verdeok}{mantém} & \textcolor{vermelhoerro}{perde} \\
        \cod{rememberSaveable} & \textcolor{verdeok}{mantém} & \textcolor{verdeok}{mantém} \\
        ViewModel & \textcolor{verdeok}{mantém} & \textcolor{verdeok}{mantém} \\
        \bottomrule
      \end{tabularx}
      \begin{atencao}
        \cod{rememberSaveable} só aceita tipos que cabem num \cod{Bundle} (primitivos, \cod{String}, \cod{Parcelable}). Para objetos maiores, use um ViewModel (Módulo 08).
      \end{atencao}
    \end{column}
  \end{columns}
\end{frame}

% =============================================================================
\section{State hoisting}

\begin{frame}[fragile]{Elevando o estado (\emph{state hoisting})}
  \begin{columns}[T]
    \begin{column}{0.52\textwidth}
\begin{codigo}[title=Componente sem estado (stateless)]{Kotlin}
@Composable
fun CampoLimite(
  valor: String,                   // estado desce
  onValorChange: (String) -> Unit, // evento sobe
  modifier: Modifier = Modifier
) {
  OutlinedTextField(
    value = valor,
    onValueChange = onValorChange,
    label = { Text("Limite (°C)") },
    modifier = modifier
  )
}

// Quem usa decide onde o estado mora
@Composable
fun TelaConfig() {
  var limite by rememberSaveable { mutableStateOf("30") }
  CampoLimite(limite, onValorChange = { limite = it })
}
\end{codigo}
    \end{column}
    \begin{column}{0.46\textwidth}
      \centering
      \begin{tikzpicture}[node distance=12mm]
        \node[caixaescura,minimum width=3.6cm] (p) {\cod{TelaConfig}\\ \tiny dona do estado};
        \node[caixa,minimum width=3.6cm,below=of p] (f) {\cod{CampoLimite}\\ \tiny stateless};
        \draw[seta] ([xshift=-8mm]p.south) -- node[rotulo,left]{estado} ([xshift=-8mm]f.north);
        \draw[setaambar] ([xshift=8mm]f.north) -- node[rotulo,right]{eventos} ([xshift=8mm]p.south);
      \end{tikzpicture}
      \begin{conceito}[Fluxo unidirecional de dados (UDF)]
        O estado desce como parâmetro e os eventos sobem como lambdas. Há uma única \textbf{fonte da verdade}.
      \end{conceito}
      \begin{dica}[Vantagens]
        Componente reutilizável, fácil de testar e de visualizar em \cod{@Preview}.
      \end{dica}
    \end{column}
  \end{columns}
\end{frame}

\begin{frame}{Onde o estado deve morar?}
  \begin{center}
  \begin{tikzpicture}[node distance=5mm and 8mm,
      decisao/.style={caixaambar,diamond,aspect=2.6,inner sep=1pt,text width=2.5cm,font=\tiny}]
    \node[decisao] (q1) {É só aparência da UI? (expandido, rolagem, foco)};
    \node[caixaescura,right=12mm of q1,text width=3.2cm] (r1) {No próprio composable,\\ com \cod{remember(Saveable)}};
    \node[decisao,below=3mm of q1] (q2) {Vários composables usam?};
    \node[caixa,right=12mm of q2,text width=3.2cm] (r2) {No ancestral comum mais baixo (hoisting)};
    \node[decisao,below=3mm of q2] (q3) {Envolve regras de negócio ou dados?};
    \node[caixaambarforte,right=12mm of q3,text width=3.2cm] (r3) {No \textbf{ViewModel} (\emph{state holder} da tela)};
    \draw[seta] (q1) -- node[rotulo,above]{sim} (r1);
    \draw[seta] (q1) -- node[rotulo,left]{não} (q2);
    \draw[seta] (q2) -- node[rotulo,above]{sim} (r2);
    \draw[seta] (q2) -- node[rotulo,left]{não} (q3);
    \draw[seta] (q3) -- node[rotulo,above]{sim} (r3);
  \end{tikzpicture}
  \end{center}
  \begin{dica}
    Regra prática: eleve o estado só até onde for necessário, nem mais nem menos.
  \end{dica}
\end{frame}

% =============================================================================
\section{Efeitos colaterais}

\begin{frame}[fragile]{APIs de efeitos}
  \begin{columns}[T]
    \begin{column}{0.5\textwidth}
\begin{codigo}[title=LaunchedEffect e DisposableEffect]{Kotlin}
@Composable
fun Monitor(bancadaId: Int, vm: PainelViewModel) {
  // roda quando entra na tela e quando a chave muda
  LaunchedEffect(bancadaId) {
    vm.carregar(bancadaId)   // coroutine
  }

  val lifecycle = LocalLifecycleOwner.current.lifecycle
  DisposableEffect(lifecycle) {
    val obs = LifecycleEventObserver { _, e ->
      if (e == Lifecycle.Event.ON_RESUME) vm.reconectar()
    }
    lifecycle.addObserver(obs)
    onDispose { lifecycle.removeObserver(obs) }
  }
}
\end{codigo}
    \end{column}
    \begin{column}{0.48\textwidth}
      \scriptsize
      \begin{tabularx}{\linewidth}{@{}l L@{}}
        \toprule
        \cabecalho{API} & \cabecalho{Quando usar} \\
        \midrule
        \cod{LaunchedEffect(chave)} & rodar \cod{suspend} ao entrar ou quando a chave muda \\
        \cod{rememberCoroutineScope} & iniciar coroutine a partir de um clique \\
        \cod{DisposableEffect} & registrar e remover \emph{listeners} \\
        \cod{SideEffect} & sincronizar com código não-Compose a cada recomposição \\
        \cod{rememberUpdatedState} & usar o valor mais recente dentro de um efeito longo \\
        \cod{produceState} & converter fonte externa em estado \\
        \cod{derivedStateOf} & estado calculado que muda menos que a origem \\
        \bottomrule
      \end{tabularx}
      \begin{atencao}
        \cod{LaunchedEffect(Unit)} roda uma única vez enquanto o composable estiver na tela.
      \end{atencao}
    \end{column}
  \end{columns}
\end{frame}

% =============================================================================
\section{Ciclo de vida}

\begin{frame}{Ciclo de vida da Activity}
  \begin{columns}[c]
    \begin{column}{0.56\textwidth}
      \centering
      \begin{tikzpicture}[node distance=2.2mm,cb/.style={caixa,minimum width=2.3cm,minimum height=5.5mm,font=\tiny}]
        \node[caixaambarforte,minimum width=2.3cm] (l) {Lançada};
        \node[cb,below=of l] (c) {\cod{onCreate()}};
        \node[cb,below=of c] (s) {\cod{onStart()}};
        \node[cb,below=of s] (r) {\cod{onResume()}};
        \node[caixaescura,minimum width=2.3cm,below=of r] (run) {Em execução};
        \node[cb,below=of run] (p) {\cod{onPause()}};
        \node[cb,below=of p] (st) {\cod{onStop()}};
        \node[cb,below=of st] (d) {\cod{onDestroy()}};
        \foreach \a/\b in {l/c,c/s,s/r,r/run,run/p,p/st,st/d}{\draw[seta] (\a) -- (\b);}
        \node[cb,right=7mm of s] (rs) {\cod{onRestart()}};
        \draw[setaambar] (st.east) -| (rs.south);
        \draw[setaambar] (rs.west) -- (s.east);
        \draw[setaambar] (p.west) -- ++(-6mm,0) |- (r.west);
        \node[rotulo,anchor=east] at ([xshift=-7mm]$(p.west)!0.5!(r.west)$) {volta ao\\ primeiro plano};
        \node[rotulo,anchor=north] at ([yshift=-1mm]rs.south east) {volta do\\ segundo plano};
      \end{tikzpicture}
    \end{column}
    \begin{column}{0.42\textwidth}
      \scriptsize
      \begin{tabularx}{\linewidth}{@{}l L@{}}
        \toprule
        \cabecalho{Estado} & \cabecalho{Significado} \\
        \midrule
        \textbf{Created} & criada, ainda invisível \\
        \textbf{Started} & visível \\
        \textbf{Resumed} & em primeiro plano, recebe toques \\
        \textbf{Paused} & parcialmente encoberta \\
        \textbf{Stopped} & invisível (segundo plano) \\
        \bottomrule
      \end{tabularx}
      \begin{atencao}[Mudanças de configuração]
        Girar a tela, trocar idioma ou tema \textbf{destrói e recria} a Activity. Por isso o estado precisa estar em \cod{rememberSaveable} ou no ViewModel.
      \end{atencao}
      \begin{atencao}[Morte do processo]
        Em segundo plano, o sistema pode encerrar o app para liberar memória. Ao voltar, só o estado salvo sobrevive.
      \end{atencao}
    \end{column}
  \end{columns}
\end{frame}

\begin{frame}[fragile]{Coletando fluxos com segurança}
  \begin{columns}[T]
    \begin{column}{0.54\textwidth}
\begin{codigo}[title=TelaPainel.kt]{Kotlin}
@Composable
fun TelaPainel(vm: PainelViewModel = viewModel()) {
  // para de coletar quando o app vai para o
  // segundo plano (abaixo de STARTED)
  val estado by vm.estado.collectAsStateWithLifecycle()

  when (val e = estado) {
    PainelUiState.Carregando -> CircularProgressIndicator()
    is PainelUiState.Erro -> Text(e.mensagem)
    is PainelUiState.Pronto -> Painel(
      leituras = e.leituras,
      onAtualizar = vm::atualizar
    )
  }
}
\end{codigo}
    \end{column}
    \begin{column}{0.44\textwidth}
      \begin{conceito}[\cod{collectAsStateWithLifecycle}]
        Converte \cod{StateFlow} em \cod{State} do Compose e \textbf{pausa a coleta} quando a tela não está visível, economizando bateria, rede e sensores.
      \end{conceito}
      \begin{dica}[Dependência]
        Vem do artefato\\ \cod{androidx.lifecycle:}\\\cod{lifecycle-runtime-compose}.
      \end{dica}
      \begin{erro}
        \cod{collectAsState()} continua coletando em segundo plano. No Android, prefira a versão \emph{WithLifecycle}.
      \end{erro}
    \end{column}
  \end{columns}
\end{frame}

% =============================================================================
\section{Navegação com Navigation 3}

\begin{frame}{Conceitos de navegação}
  \begin{columns}[c]
    \begin{column}{0.5\textwidth}
      \small
      \begin{itemize}
        \item Apps Compose costumam ter \textbf{uma Activity} e várias telas (\emph{destinos}).
        \item A \termo{pilha de navegação} (\emph{back stack}) guarda o histórico; o botão/gesto voltar remove o topo.
        \item \textbf{Navigation 3} (Jetpack, estável desde nov./2025) representa a pilha como uma \textbf{lista observável de chaves} controlada pelo próprio app.
      \end{itemize}
      \begin{conceito}[Peças]
        \cod{NavKey} (chave de cada tela) $\cdot$ \cod{rememberNavBackStack} (a pilha) $\cdot$ \cod{NavDisplay} (exibe o topo) $\cdot$ \cod{entryProvider} (chave $\rightarrow$ tela).
      \end{conceito}
    \end{column}
    \begin{column}{0.48\textwidth}
      \centering
      \begin{tikzpicture}[node distance=1mm]
        \node[caixa,minimum width=3.2cm] (a) {\cod{Lista}};
        \node[caixa,minimum width=3.2cm,above=of a] (b) {\cod{Detalhe(id = 42)}};
        \node[caixaambarforte,minimum width=3.2cm,above=of b] (c) {\cod{Reserva(id = 42)}};
        \node[rotulo,right=2mm of c] {topo: visível};
        \draw[setaambar] ([xshift=-4mm]c.north west) -- ++(0,4mm) node[above,rotulo]{\cod{add()}};
        \draw[seta] ([xshift=4mm]c.north east) ++(0,4mm) -- ([xshift=4mm]c.north east) ;
        \node[rotulo] at ([xshift=12mm,yshift=6mm]c.north east) {\cod{removeLastOrNull()}};
        \node[pontilhado,fit=(a)(b)(c)] {};
      \end{tikzpicture}
      \begin{dica}
        Muitos projetos existentes usam \textbf{Navigation Compose} (Nav2, com \cod{NavHost} e \cod{NavController}). Os conceitos são os mesmos; ao final mostramos a equivalência.
      \end{dica}
    \end{column}
  \end{columns}
\end{frame}

\begin{frame}[fragile]{Navigation 3: chaves e pilha}
  \begin{columns}[T]
    \begin{column}{0.5\textwidth}
\begin{codigo}[title=app/build.gradle.kts]{Kotlin}
plugins {
  alias(libs.plugins.kotlin.serialization)
}
dependencies {
  implementation(libs.androidx.navigation3.runtime)
  implementation(libs.androidx.navigation3.ui)
  implementation(
    libs.androidx.lifecycle.viewmodel.navigation3)
  implementation(libs.kotlinx.serialization.core)
}
\end{codigo}
\begin{codigo}[title=Rotas.kt]{Kotlin}
@Serializable data object Lista : NavKey
@Serializable data class Detalhe(val id: Int) : NavKey
@Serializable data class Reserva(val id: Int) : NavKey
\end{codigo}
    \end{column}
    \begin{column}{0.48\textwidth}
      \begin{conceito}[Chaves tipadas]
        Cada tela é uma classe \cod{@Serializable} que implementa \cod{NavKey}. Os argumentos são propriedades: sem strings de rota e sem erros de digitação.
      \end{conceito}
      \begin{conceito}[Pilha persistente]
        \cod{rememberNavBackStack(Lista)} cria a pilha com a tela inicial e a salva automaticamente em rotações e na morte do processo, por isso as chaves são serializáveis.
      \end{conceito}
      \begin{dica}
        Passe só identificadores (\cod{id}) nas chaves; a tela de destino busca os dados.
      \end{dica}
    \end{column}
  \end{columns}
\end{frame}

\begin{frame}[fragile]{Navigation 3: \cod{NavDisplay}}
  \begin{columns}[T]
    \begin{column}{0.6\textwidth}
\begin{codigo}[title=NavegacaoApp.kt]{Kotlin}
@Composable
fun NavegacaoApp() {
  val pilha = rememberNavBackStack(Lista)
  NavDisplay(
    backStack = pilha,
    onBack = { pilha.removeLastOrNull() },
    entryDecorators = listOf(
      rememberSaveableStateHolderNavEntryDecorator(),
      rememberViewModelStoreNavEntryDecorator()), // VM por tela
    entryProvider = entryProvider {
      entry<Lista> {
        TelaLista(onAbrir = { id -> pilha.add(Detalhe(id)) })
      }
      entry<Detalhe> { chave ->
        TelaDetalhe(chave.id,
          onReservar = { pilha.add(Reserva(chave.id)) })
      }
      entry<Reserva> { chave -> TelaReserva(chave.id) }
    }
  )
}
\end{codigo}
    \end{column}
    \begin{column}{0.38\textwidth}
      \small
      \begin{itemize}
        \item \textbf{Navegar} = \cod{pilha.add(chave)}.
        \item \textbf{Voltar} = \cod{pilha.removeLastOrNull()}.
        \item O gesto de voltar do sistema chama \cod{onBack}.
        \item Os decoradores salvam o estado de cada tela e criam um ViewModel por entrada da pilha.
      \end{itemize}
      \begin{atencao}[Telas não conhecem a pilha]
        \cod{TelaLista} recebe \cod{onAbrir} como lambda, não a pilha. Assim ela continua reutilizável e testável.
      \end{atencao}
    \end{column}
  \end{columns}
\end{frame}

\begin{frame}[fragile]{Navegação principal com barra inferior}
  \begin{columns}[T]
    \begin{column}{0.56\textwidth}
\begin{codigo}[title=Abas principais]{Kotlin}
enum class Aba(val chave: NavKey, val rotulo: String,
               val icone: ImageVector) {
  PAINEL(Painel, "Painel", Icons.Default.Dashboard),
  EQUIP(Lista, "Equipamentos", Icons.Default.Memory),
  CONFIG(Config, "Ajustes", Icons.Default.Settings)
}

NavigationBar {
  Aba.entries.forEach { aba ->
    NavigationBarItem(
      selected = pilha.first() == aba.chave,
      onClick = {           // troca a raiz da pilha
        pilha.clear(); pilha.add(aba.chave)
      },
      icon = { Icon(aba.icone, null) },
      label = { Text(aba.rotulo) }
    )
  }
}
\end{codigo}
    \end{column}
    \begin{column}{0.42\textwidth}
      \begin{dica}[Uma pilha por aba]
        Em apps maiores, mantenha uma pilha para cada aba e exiba a da aba ativa. Assim cada aba lembra até onde o usuário navegou. Os \emph{recipes} oficiais do Navigation 3 trazem esse padrão pronto.
      \end{dica}
      \begin{conceito}[Nav2 $\leftrightarrow$ Nav3]
        \scriptsize
        \begin{tabular}{@{}l l@{}}
          \cod{NavHost} & \cod{NavDisplay} \\
          \cod{navigate(r)} & \cod{pilha.add(r)} \\
          \cod{popBackStack()} & \cod{removeLastOrNull()} \\
          \cod{composable<R>\{\}} & \cod{entry<R>\{\}} \\
          \cod{toRoute<R>()} & a própria chave \\
        \end{tabular}
      \end{conceito}
    \end{column}
  \end{columns}
\end{frame}

% =============================================================================
\section{Animações}

\begin{frame}[fragile]{Animações por valor alvo}
  \begin{columns}[T]
    \begin{column}{0.56\textwidth}
\begin{codigo}[title=CartaoMedidor.kt]{Kotlin}
@Composable
fun CartaoMedidor(temp: Double, expandido: Boolean) {
  val alerta = temp > 30
  val altura by animateDpAsState(
    if (expandido) 160.dp else 72.dp, label = "altura")
  val cor by animateColorAsState(
    if (alerta) Ambar else MaterialTheme.colorScheme.surface,
    label = "cor")
  Card(
    modifier = Modifier.height(altura),
    colors = CardDefaults.cardColors(containerColor = cor)
  ) {
    AnimatedContent(temp, label = "temp") { t ->
      Text("%.1f °C".format(t))
    }
    AnimatedVisibility(visible = alerta) {
      Text("Acima do limite!")
    }
  }
}
\end{codigo}
    \end{column}
    \begin{column}{0.42\textwidth}
      \begin{conceito}[Valor alvo]
        Você informa o valor final; o Compose anima do valor atual até ele a cada mudança de estado. Cobre a maioria dos casos.
      \end{conceito}
      \scriptsize
      \begin{tabular}{@{}l l@{}}
        \toprule
        \cabecalho{API} & \cabecalho{Anima} \\
        \midrule
        \cod{animate*AsState} & um valor (cor, \cod{Dp}, \cod{Float}) \\
        \cod{AnimatedVisibility} & entrada e saída \\
        \cod{AnimatedContent} & troca de conteúdo \\
        \cod{animateContentSize} & mudança de tamanho \\
        \bottomrule
      \end{tabular}
      \normalsize
      \begin{dica}
        O parâmetro \cod{label} identifica a animação no \emph{Animation Preview} do Android Studio.
      \end{dica}
    \end{column}
  \end{columns}
\end{frame}

\begin{frame}[fragile]{Animações contínuas, molas e transições}
  \begin{columns}[T]
    \begin{column}{0.54\textwidth}
\begin{codigo}[title=Ícone pulsando]{Kotlin}
@Composable
fun IconeAlerta() {
  val pulso = rememberInfiniteTransition(label = "pulso")
  val escala by pulso.animateFloat(
    initialValue = 1f, targetValue = 1.3f,
    animationSpec = infiniteRepeatable(
      tween(1000), RepeatMode.Reverse),
    label = "escala")
  Icon(Icons.Default.Warning, "Alerta", tint = Ambar,
    modifier = Modifier.graphicsLayer {
      scaleX = escala; scaleY = escala
    })
}
\end{codigo}
    \end{column}
    \begin{column}{0.44\textwidth}
\begin{codigo}[title=Painel com mola]{Kotlin}
val deslocamento by animateDpAsState(
  if (aberto) 0.dp else 300.dp,
  animationSpec = spring(
    dampingRatio = 0.5f,   // MediumBouncy
    stiffness = 200f),     // StiffnessLow
  label = "painel")
\end{codigo}
      \begin{conceito}[Mola \emph{versus} \cod{tween}]
        \cod{tween} tem duração e curva fixas; \cod{spring} segue a física e continua suave se o alvo mudar no meio do caminho.
      \end{conceito}
      \begin{dica}[Transições entre telas]
        \cod{SharedTransitionLayout} faz a foto ``voar'' da lista ao detalhe; o Navigation 3 anima o voltar preditivo.
      \end{dica}
    \end{column}
  \end{columns}
\end{frame}

% =============================================================================
\section{Fechamento}

\begin{frame}{Resumo}
  \begin{resumo}
    \begin{itemize}
      \item \cod{mutableStateOf} + \cod{remember} criam estado observável; mudar o estado dispara a recomposição.
      \item \cod{rememberSaveable} sobrevive à rotação; o ViewModel guarda o estado da tela.
      \item \emph{State hoisting}: estado desce, eventos sobem, uma fonte da verdade.
      \item Efeitos colaterais só dentro de \cod{LaunchedEffect}, \cod{DisposableEffect} e afins.
      \item Rotação recria a Activity; colete fluxos com \cod{collectAsStateWithLifecycle}.
      \item Navigation 3: chaves \cod{NavKey}, pilha observável e \cod{NavDisplay}.
      \item Animações: \cod{animate*AsState}, \cod{AnimatedVisibility}, \cod{AnimatedContent} e molas.
    \end{itemize}
  \end{resumo}
\end{frame}

\begin{frame}{Glossário (1/2)}
  \begin{glossario}
    Estado & valor que muda com o tempo e afeta o que a tela mostra \\
    Estado derivado & valor calculado a partir de outros; não se guarda \\
    Fonte única da verdade & cada dado tem um único dono que pode alterá-lo \\
    Fluxo unidirecional & o estado desce até a interface; os eventos sobem até o dono \\
    \emph{State hoisting} & mover o estado para quem chama, deixando o componente sem estado \\
    \cod{remember} & guarda um valor entre recomposições \\
    \cod{rememberSaveable} & guarda também na rotação e na morte do processo \\
    Efeito colateral & ação fora da descrição da tela: rede, registro, temporizador \\
    \cod{LaunchedEffect} & executa um efeito suspenso ligado à composição \\
  \end{glossario}
\end{frame}

\begin{frame}{Glossário (2/2)}
  \begin{glossario}
    Ciclo de vida & estados da \cod{Activity}: \emph{Created}, \emph{Started}, \emph{Resumed}\ldots \\
    Morte do processo & encerramento do app em segundo plano para liberar memória \\
    Pilha de retorno & histórico LIFO dos destinos visitados \\
    \cod{NavKey} & chave que identifica um destino no Navigation 3 \\
    Link profundo & endereço que abre o app direto em um destino \\
    Voltar preditivo & gesto que mostra a tela anterior antes de concluir o voltar \\
    Mola (\cod{spring}) & animação definida por rigidez e amortecimento, não por duração \\
  \end{glossario}
\end{frame}

\begin{frame}{Questões de revisão}
  \begin{questoes}
    \item Por que uma variável comum dentro de um composable não serve para guardar estado?
    \item Explique o fluxo unidirecional de dados com o exemplo de um campo de busca.
    \item Modele a tela de equipamentos como máquina de estados. Por que uma \cod{sealed interface} e não três booleanos?
    \item O que sobrevive à rotação e à morte do processo com \cod{remember}, \cod{rememberSaveable} e \cod{ViewModel}?
    \item O que deve acontecer com a leitura de sensores quando o app vai para o segundo plano? Por quê?
    \item Desenhe a pilha de retorno depois de Painel, Equipamentos, Detalhe e um voltar.
    \item Como um link profundo para a reserva 42 deve montar a pilha? Que validações fazer?
    \item Por que uma animação com \cod{spring} lida melhor com interrupções que um \cod{tween}?
  \end{questoes}
\end{frame}

\begin{frame}{Exercícios}
  \begin{exercicio}[Exercício 1 --- Rotação]
    Crie um formulário com três campos, um usando \cod{remember} e dois com \cod{rememberSaveable}. Gire o emulador e explique o que aconteceu.
  \end{exercicio}
  \begin{exercicio}[Exercício 2 --- Hoisting]
    Transforme um \cod{SeletorDeFaixa} com estado interno (mínimo e máximo) em um componente \emph{stateless} com \cod{RangeSlider}. Escreva dois \cod{@Preview}s com valores diferentes.
  \end{exercicio}
  \begin{exercicio}[Exercício 3 --- Ciclo de vida]
    Registre um \cod{Log} em cada callback do ciclo de vida da \cod{MainActivity}. Gire a tela, vá para a tela inicial e volte. Desenhe a sequência observada.
  \end{exercicio}
  \begin{projeto}[LabMonitor --- Etapa 7]
    Implemente a navegação: \textbf{Painel}, \textbf{Equipamentos} (lista $\rightarrow$ detalhe $\rightarrow$ reserva) e \textbf{Ajustes} com \cod{NavigationBar}. O formulário de reserva deve sobreviver à rotação.
  \end{projeto}
\end{frame}

\begin{frame}{Para aprofundar}
  \small
  \begin{itemize}
    \item Estado e Jetpack Compose: \url{https://developer.android.com/develop/ui/compose/state}
    \item Onde elevar o estado: \url{https://developer.android.com/develop/ui/compose/state-hoisting}
    \item Efeitos colaterais: \url{https://developer.android.com/develop/ui/compose/side-effects}
    \item Ciclo de vida da Activity: \url{https://developer.android.com/guide/components/activities/activity-lifecycle}
    \item Navigation 3: \url{https://developer.android.com/guide/navigation/navigation-3}
    \item Recipes do Navigation 3: \url{https://github.com/android/nav3-recipes}
  \end{itemize}
\end{frame}

\begin{frame}{Respostas das questões de revisão}
  \begin{respostas}
    \item Ela volta ao valor inicial a cada recomposição, e alterá-la não avisa o Compose. Use \cod{mutableStateOf} dentro de \cod{remember}.
    \item O \cod{TextField} mostra \cod{busca} (estado que desce) e chama \cod{onBuscaChange} (evento que sobe); o \cod{ViewModel} atualiza o estado, que volta à tela.
    \item Estados \cod{Carregando}, \cod{Sucesso(itens)}, \cod{Vazio} e \cod{Erro(msg)}. A classe selada só admite esses quatro e o \cod{when} obriga a tratar todos; booleanos permitem combinações impossíveis.
    \item \cod{remember}: nenhuma das duas. \cod{rememberSaveable}: ambas. \cod{ViewModel}: rotação sim; morte do processo só com \cod{SavedStateHandle}.
    \item Deve parar, porque gasta bateria e ninguém está vendo. \cod{collectAsStateWithLifecycle} para de coletar quando a tela sai de \emph{Started}.
    \item Antes do voltar: Painel, Equipamentos, Detalhe (topo). Depois: Painel, Equipamentos (topo).
    \item Com uma pilha sintética Painel, Reserva(42), para o voltar levar ao Painel. Validar se a reserva existe e se o usuário pode vê-la.
    \item A mola parte da posição e da velocidade atuais e segue a física até o novo alvo; o \cod{tween} recomeça com duração fixa e muda a velocidade de forma brusca.
  \end{respostas}
\end{frame}

\encerramento{Módulo 08 --- Arquitetura de apps: MVVM, camadas, repositório, casos de uso e injeção de dependência com Hilt.}

\end{document}
